% AI review of the final deck: reviewer setup with the review prompt on one
% frame, then the grade with the iteration history.
% Grade macros (\resReviewScore, \resReviewVerbal) come from review.tex,
% maintained by the orchestrator.
% Grade of the AI presentation review of the final deck, copied verbatim by
% the orchestrator from presentation-reviewer: \resReviewScore and
% \resReviewVerbal from the final iteration (4), the other scores from
% iterations 1 to 3. Data for slides 19 and 20; do not edit on the slide side.
\newcommand{\resReviewFirstScore}{7}
\newcommand{\resReviewSecondScore}{8}
\newcommand{\resReviewThirdScore}{8}
\newcommand{\resReviewScore}{8}
\newcommand{\resReviewVerbal}{Prezentacja spełnia wszystkie warunki obowiązkowe, a po skróceniu części eksperymentalnej nadal obejmuje każdy oceniany element, czyli regułę z uzasadnieniem, wynik globalny z porównaniem oraz przykład lokalny z kryterium wyboru. Wykresy są czytelne, liczby pochodzą z programu, a długość talii mieści się w zalecanym zakresie. Do poprawy pozostają opis narzędzi recenzentów, który nie zgadza się z dziennikiem pracy, oraz słabe odniesienie przykładu lokalnego do całego zbioru.}


\begin{frame}[fragile]
  \frametitle{Recenzent ocenia talię według trzech źródeł kryteriów z promptu}
  \begin{itemize}
    \item Agent presentation-reviewer czyta skompilowany PDF i źródła LaTeX.
      Definicja roli daje mu tylko narzędzia odczytu, a przy tej talii
      pilnowało tego polecenie.
    \item Każda iteracja recenzji daje znaleziska z wagami oraz ocenę słowną
      i liczbową. Ocena jest doradcza.
  \end{itemize}
  {\footnotesize
\begin{verbatim}
[...]
Review slides/final/main.pdf and its .tex sources
against: (1) [...]zasady_zaliczenia_projektu_v1.pdf, mandatory
conditions and slide groups; (2) [...]Projekt_1.pdf
p. 6-9, final deck content, slides A-F, experiment 5.2,
table 2; (3) docs/presentation-style.md, assertion-evidence
titles, budgets, palette, banned elements. Report findings
with severity and a suggested fix. End with a verbal grade
and a score 1-10.
[...]
\end{verbatim}}
\end{frame}

\begin{frame}
  \frametitle{Recenzent AI ocenił prezentację końcową na \resReviewScore{}\,/\,10}
  „\resReviewVerbal''
  \vskip3mm
  Kolejne iteracje dały \resReviewFirstScore{}\,/\,10,
  \resReviewSecondScore{}\,/\,10, \resReviewThirdScore{}\,/\,10
  i \resReviewScore{}\,/\,10. Porównanie z dziennikiem i danymi wykryło
  nieścisłości, których kompilacja nie pokazuje.
  \vskip3mm
  {\footnotesize\color{mnMuted}
    Ocena pochodzi z końcowej iteracji recenzji. Przebieg iteracji
    dokumentuje dziennik pracy z AI.\par}
\end{frame}
